\documentclass[../main]{subfiles}
\ifSubfilesClassLoaded{\setcounter{chapter}{3}}{}
\begin{document}

\chapter{Tabel Kebenaran, Tautologi, Kontradiksi, dan Kontingensi}

\begin{subcpmk}
  \item \textbf{Sub-CPMK 1.2:} Mahasiswa mampu menggunakan tabel kebenaran untuk mengevaluasi pernyataan majemuk serta mengidentifikasi tautologi, kontradiksi, dan kontingensi.
\end{subcpmk}

% ============================================================
% MATERI POKOK
% ============================================================
\section{Tautologi, Kontradiksi, dan Kontingensi}

Klasifikasi proposisi majemuk berdasarkan nilai kebenarannya merupakan fondasi penting sebelum membahas ekuivalensi logis dan inferensi formal \cite{rosen2019}. Secara umum, proposisi majemuk dikelompokkan menjadi tiga jenis, yaitu tautologi, kontradiksi, dan kontingensi.

\subsection{Tautologi}

\begin{definisi}[Tautologi]
\logic{Tautologi} adalah proposisi majemuk yang \textbf{selalu bernilai benar (True)} untuk setiap kombinasi nilai kebenaran proposisi atomik penyusunnya.
\end{definisi}

Dalam konteks pemrograman atau basis data, ekspresi yang selalu benar tidak memberikan penyaringan kondisi yang efektif karena hasil evaluasinya selalu meloloskan.

\begin{contoh}
Contoh tautologi paling sederhana adalah \textbf{hukum tiada jalan tengah}: $p \vee \neg p$. \\
Pernyataan ini selalu bernilai benar, apa pun nilai kebenaran $p$.
\end{contoh}

\subsection{Kontradiksi}

\begin{definisi}[Kontradiksi]
\logic{Kontradiksi} adalah proposisi majemuk yang \textbf{selalu bernilai salah (False)} untuk setiap kombinasi nilai kebenaran proposisi atomik penyusunnya.
\end{definisi}

Dalam penerapan komputasi, kondisi kontradiktif dapat menyebabkan cabang program tidak pernah dieksekusi (\textit{dead code}).

\begin{contoh}
Contoh kontradiksi paling sederhana adalah \textbf{hukum non-kontradiksi}: $p \wedge \neg p$. \\
Pernyataan ini selalu bernilai salah karena $p$ dan $\neg p$ tidak mungkin benar pada saat yang sama.
\end{contoh}

\subsection{Kontingensi}

\begin{definisi}[Kontingensi]
\logic{Kontingensi} adalah proposisi majemuk yang nilai kebenarannya \textbf{bergantung} pada kombinasi nilai proposisi atomik penyusunnya; pada sebagian kombinasi bernilai benar, dan pada kombinasi lain bernilai salah.
\end{definisi}

Sebagian besar ekspresi logika dalam praktik, seperti $p \rightarrow q$ atau $(p \vee q) \wedge \neg r$, termasuk kontingensi. Oleh karena itu, evaluasi umumnya memerlukan tabel kebenaran atau manipulasi aljabar yang tepat.

\section{Konsep Ekuivalensi Logis Melalui Tabel Kebenaran}

Dua ekspresi logika dapat tampak berbeda secara bentuk, tetapi tetap memiliki makna logis yang sama. Dalam konteks matematika maupun rekayasa perangkat lunak, pengujian kesetaraan tersebut disebut uji ekuivalensi logis \cite{rosen2019}.

\subsection{Parameter Ekuivalensi Mutlak}

\begin{definisi}[Ekuivalen Logis]
Dua proposisi majemuk $P$ dan $Q$ disebut \logic{ekuivalen logis} (ditulis $\equiv$) apabila dan hanya apabila kolom nilai kebenaran keduanya identik untuk seluruh kombinasi nilai variabel penyusun. \\
Secara aljabar, hal ini setara dengan menunjukkan bahwa $(P \leftrightarrow Q)$ merupakan \textbf{tautologi}.
\end{definisi}

\subsection{Verifikasi Ekuivalensi dengan Tabel Kebenaran}

Tabel kebenaran merupakan metode langsung untuk memverifikasi ekuivalensi. Meskipun bersifat enumeratif, metode ini memberikan hasil yang pasti. Perhatikan pengujian berikut untuk salah satu bentuk Hukum De Morgan:
\begin{center}
``Apakah $\neg(p \wedge q)$ ekuivalen dengan $\neg p \vee \neg q$?''
\end{center}

Berikut tabel kebenaran untuk variabel $p$ dan $q$:

\begin{table}[H]
\centering
\begin{tabular}{|c|c||c|c||c|c||c|}
\hline
\textbf{$p$} & \textbf{$q$} & \textbf{$\neg p$} & \textbf{$\neg q$} & \textbf{Langkah $p \wedge q$} & \textbf{Objek A: $\neg(p \wedge q)$} & \textbf{Objek B: $\neg p \vee \neg q$} \\
\hline
T & T & F & F & T & \textbf{F} & \textbf{F} \\
T & F & F & T & F & \textbf{T} & \textbf{T} \\
F & T & T & F & F & \textbf{T} & \textbf{T} \\
F & F & T & T & F & \textbf{T} & \textbf{T} \\
\hline
\end{tabular}
\end{table}

\textbf{Analisis hasil:} \\
Kolom \textbf{Objek A} dan \textbf{Objek B} menghasilkan pola nilai yang sama, yaitu $(F, T, T, T)$. Karena nilainya identik pada seluruh baris, dapat disimpulkan bahwa:
\[
\neg(p \wedge q) \equiv \neg p \vee \neg q
\]

Untuk ekspresi dengan banyak variabel, jumlah baris tabel kebenaran meningkat secara eksponensial. Oleh sebab itu, selain tabel kebenaran, digunakan pula hukum-hukum aljabar logika untuk pembuktian yang lebih efisien \cite{wiki_proplogic}.

\section{Hukum-hukum Dasar Aljabar Proposisi}

Hukum-hukum aljabar proposisi digunakan untuk menyederhanakan ekspresi logika tanpa harus selalu membuat tabel kebenaran lengkap. Hukum-hukum ini telah tervalidasi secara formal dan banyak dipakai dalam analisis logika, pemrograman, serta perancangan rangkaian digital \cite{rosen2019}.

Secara praktis, hukum-hukum tersebut dapat dikelompokkan menurut fungsinya.

\subsection{Hukum Pengaruh Absolut (Dominasi \& Identitas)}
Nilai konstan benar (T) atau salah (F) dapat memengaruhi hasil ekspresi secara langsung melalui hukum dominasi dan identitas.

\begin{itemize}
    \item \textbf{Hukum Dominasi (Dominasi Mutlak / Null Laws)} \\
    Pada disjungsi dengan T, hasil selalu T. Pada konjungsi dengan F, hasil selalu F.
    \begin{align*}
        p \vee \mathbf{T} &\equiv \mathbf{T} \\
        p \wedge \mathbf{F} &\equiv \mathbf{F}
    \end{align*}

    \item \textbf{Hukum Identitas (Hukum Cermin Transparan)} \\
    Pada konjungsi dengan T atau disjungsi dengan F, nilai variabel tetap.
    \begin{align*}
        p \wedge \mathbf{T} &\equiv p \\
        p \vee \mathbf{F} &\equiv p
    \end{align*}
\end{itemize}

\subsection{Hukum Konflik Internal (Negasi \& Komplemen)}
Bagian ini menjelaskan relasi variabel dengan negasinya.

\begin{itemize}
    \item \textbf{Hukum Negasi (Komplemen Akhir)} \\
    Kombinasi variabel dengan negasinya menghasilkan tautologi atau kontradiksi.
    \begin{align*}
        p \vee \neg p &\equiv \mathbf{T} \quad \text{(Hukum Tiada Jalan Tengah)} \\
        p \wedge \neg p &\equiv \mathbf{F} \quad \text{(Hukum Non-Kontradiksi)}
    \end{align*}

    \item \textbf{Hukum Involusi (Pembatalan Negasi Ganda)} \\
    Negasi ganda mengembalikan nilai semula.
    \begin{align*}
        \neg(\neg p) &\equiv p 
    \end{align*}
\end{itemize}

\subsection{Hukum Redundansi Berlebihan (Idempoten \& Absorpsi)}
Hukum ini dipakai untuk menghilangkan pengulangan yang tidak diperlukan.

\begin{itemize}
    \item \textbf{Hukum Idempoten (Sia-sia Kembar Identik)} \\
    Pengulangan variabel yang sama pada operator sejenis tidak mengubah nilai.
    \begin{align*}
        p \vee p &\equiv p \\
        p \wedge p &\equiv p
    \end{align*}

    \item \textbf{Hukum Absorpsi (Penyerapan Raksasa Hitam)} \\
    Variabel utama dapat menyerap bentuk majemuk yang memuat variabel itu sendiri.
    \begin{align*}
        p \vee (p \wedge q) &\equiv p \\
        p \wedge (p \vee q) &\equiv p
    \end{align*}
\end{itemize}
Hukum-hukum di atas sangat membantu dalam penyederhanaan ekspresi logika secara sistematis \cite{wiki_proplogic}.

\section{Hukum Lanjutan: Pertukaran, Distributif, dan De Morgan}

Setelah memahami hukum dasar, langkah berikutnya adalah menggunakan hukum lanjutan untuk mengubah susunan ekspresi tanpa mengubah maknanya. Hukum yang utama pada bagian ini adalah komutatif, asosiatif, distributif, dan De Morgan \cite{rosen2019}.

\subsection{Sifat Komutatif dan Asosiatif}

Sebagaimana pada aljabar aritmetika, operator AND dan OR memiliki sifat pertukaran urutan dan pengelompokan tertentu.

\begin{itemize}
    \item \textbf{Hukum Komutatif} \\
    Urutan operand dapat ditukar tanpa mengubah hasil.
    \begin{align*}
        p \vee q &\equiv q \vee p \\
        p \wedge q &\equiv q \wedge p 
    \end{align*}

    \item \textbf{Hukum Asosiatif} \\
    Pengelompokan operand pada operator sejenis dapat diubah tanpa memengaruhi nilai akhir.
    \begin{align*}
        (p \vee q) \vee r &\equiv p \vee (q \vee r) \\
        (p \wedge q) \wedge r &\equiv p \wedge (q \wedge r) 
    \end{align*}
\end{itemize}

\subsection{Hukum Distributif}

Hukum distributif memungkinkan operator di luar kurung didistribusikan ke masing-masing komponen di dalam kurung, seperti analogi distributif pada aljabar aritmetika.

\begin{align*}
    p \vee (q \wedge r) &\equiv (p \vee q) \wedge (p \vee r) \\
    p \wedge (q \vee r) &\equiv (p \wedge q) \vee (p \wedge r) 
\end{align*}
Kedua bentuk distributif di atas dapat digunakan dua arah: untuk mengembangkan ekspresi maupun untuk melakukan faktorisasi.

\subsection{Hukum De Morgan}

Hukum De Morgan menjelaskan bagaimana negasi bekerja pada ekspresi majemuk yang mengandung operator AND atau OR \cite{wiki_proplogic}.

\begin{teorema}[Hukum De Morgan]
Negasi yang masuk ke dalam kurung mengubah setiap komponen menjadi negasinya dan menukar operator utama: AND menjadi OR, serta OR menjadi AND.
\end{teorema}

\begin{align*}
    \neg(p \wedge q) &\equiv \neg p \vee \neg q \\
    \neg(p \vee q) &\equiv \neg p \wedge \neg q
\end{align*}

Penguasaan hukum distributif dan De Morgan sangat penting untuk penyederhanaan ekspresi logika, baik pada analisis matematis maupun perancangan sirkuit digital.

\section{Transformasi Ekuivalensi Bersyarat (Implikasi \& Biimplikasi)}

Ekspresi logika sering menggunakan operator implikasi ($\rightarrow$) dan biimplikasi ($\leftrightarrow$). Untuk menyederhanakannya, kedua bentuk tersebut dapat ditransformasikan ke operator dasar AND, OR, dan NOT \cite{rosen2019}.

\subsection{Transformasi Implikasi}

Implikasi $(p \rightarrow q)$ merepresentasikan pernyataan ``Jika $p$, maka $q$''. Bentuk ini mempunyai ekuivalensi standar berikut:

\begin{teorema}[Ekuivalensi Implikasi]
Struktur $(p \rightarrow q)$ ekuivalen dengan bentuk disjungsi $(\neg p \vee q)$.
\end{teorema}

Bentuk yang sama juga dapat dituliskan sebagai:
\[
p \rightarrow q \equiv \neg (p \wedge \neg q)
\]

\textbf{Contoh singkat:} \\
``Jika Anda menyumbang darah ($p$), maka saya memberi apresiasi ($q$).'' \\
Secara logika, pernyataan tersebut ekuivalen dengan ``Anda tidak menyumbang darah ($\neg p$), atau saya memberi apresiasi ($q$)'', yaitu $\neg p \vee q$.

\subsection{Transformasi Biimplikasi}

Biimplikasi menyatakan hubungan dua arah: ``jika dan hanya jika''.

\begin{teorema}[Ekuivalensi Biimplikasi]
Struktur $(p \leftrightarrow q)$ dapat diuraikan sebagai:
\begin{align*}
    p \leftrightarrow q &\equiv (p \rightarrow q) \wedge (q \rightarrow p) \\
    \text{dengan substitusi implikasi: } & \\
    p \leftrightarrow q &\equiv (\neg p \vee q) \wedge (\neg q \vee p)
\end{align*}
\end{teorema}
Bentuk ekuivalen lain yang lazim digunakan adalah:
\[
p \leftrightarrow q \equiv (p \wedge q) \vee (\neg p \wedge \neg q)
\]

\subsection{Pola Penyederhanaan Campuran}

Dalam praktik, terdapat bentuk campuran yang sering muncul dan dapat disederhanakan secara cepat:

\begin{align*}
    p \vee (\neg p \wedge q) &\equiv (p \vee q) \\
    p \wedge (\neg p \vee q) &\equiv (p \wedge q)
\end{align*}

Pembuktian untuk bentuk pertama:
\begin{enumerate}
    \item Terapkan distributif: $\equiv (p \vee \neg p) \wedge (p \vee q)$
    \item Terapkan hukum negasi: $\equiv \mathbf{T} \wedge (p \vee q)$
    \item Terapkan identitas: $\therefore \equiv (p \vee q)$
\end{enumerate}
Pola-pola ini bermanfaat untuk mempercepat penyederhanaan ekspresi pada analisis logika maupun refactoring kode.

\section{Strategi Penyederhanaan Ekspresi Bertahap}

Ketika jumlah variabel bertambah, tabel kebenaran menjadi semakin besar. Oleh karena itu, penyederhanaan aljabar dilakukan secara bertahap dengan memanfaatkan hukum-hukum logika \cite{wiki_proplogic,rosen2019}. Tujuannya adalah memperoleh bentuk yang lebih ringkas tanpa mengubah nilai logis ekspresi.

\subsection{Urutan Operasi Penyederhanaan Aljabar Boolean}

Salah satu urutan kerja yang umum digunakan adalah sebagai berikut:
\begin{enumerate}
    \item \textbf{Ubah implikasi/biimplikasi}: transformasikan $\rightarrow$ dan $\leftrightarrow$ ke bentuk $\wedge, \vee, \neg$.
    \item \textbf{Dorong negasi ke dalam}: gunakan Hukum De Morgan dan negasi ganda bila ada bentuk $\neg(\cdots)$.
    \item \textbf{Sederhanakan pola langsung}: terapkan identitas, dominasi, idempoten, dan absorpsi.
    \item \textbf{Gunakan distributif bila diperlukan}: lakukan ekspansi atau faktorisasi untuk membuka pola yang lebih sederhana.
\end{enumerate}

\subsection{Contoh Penyederhanaan Langkah Demi Langkah}

\begin{contoh}[Penyederhanaan Ekspresi]
Sederhanakan ekspresi berikut:
\[
\neg(p \vee \neg q) \vee (p \wedge q \wedge \neg r)
\]

\textbf{Langkah penyelesaian:}
\vspace{2mm}
    
\begin{tabular}{p{0.4cm} p{6.2cm} p{6.5cm}}
    1. & $\neg(p \vee \neg q) \vee (p \wedge (q \wedge \neg r))$ & Asosiatif untuk merapikan bagian kanan. \\ 
    2. & $(\neg p \wedge \neg(\neg q)) \vee (p \wedge (q \wedge \neg r))$ & De Morgan pada bagian kiri. \\ 
    3. & $(\neg p \wedge q) \vee (p \wedge (q \wedge \neg r))$ & Involusi, $\neg(\neg q)\equiv q$. \\ 
    4. & $(\neg p \wedge q) \vee ((p \wedge \neg r) \wedge q)$ & Komutatif pada bagian kanan. \\ 
    5. & $q \wedge ( \neg p \vee (p \wedge \neg r) )$ & Faktorisasi distributif. \\
    6. & $q \wedge ( (\neg p \vee p) \wedge (\neg p \vee \neg r) )$ & Distributif pada bagian dalam. \\
    7. & $q \wedge ( \mathbf{T} \wedge (\neg p \vee \neg r) )$ & Hukum negasi, $\neg p \vee p \equiv \mathbf{T}$. \\
    8. & $\mathbf{q \wedge (\neg p \vee \neg r)}$ & Hukum identitas.
\end{tabular} 
\vspace{2mm}

Hasil akhir yang lebih sederhana adalah $\mathbf{q \wedge (\neg p \vee \neg r)}$. Bentuk ini ekuivalen dengan ekspresi awal, tetapi lebih efisien untuk dianalisis dan diimplementasikan.
\end{contoh}

\section{Studi Kasus: Optimasi Kode Bercabang \textit{If-Else} dan Gerbang Logika}

Penerapan aljabar logika tidak berhenti pada pembuktian teoritis. Dalam praktik, aljabar logika digunakan untuk menyederhanakan percabangan program (\textit{software}) dan mengoptimalkan rancangan gerbang digital (\textit{hardware}) \cite{rosen2019}.

\subsection{Penyederhanaan Percabangan \textit{If-Else}}

Ekspresi logika yang terlalu kompleks dalam kode dapat menyulitkan pemeliharaan dan meningkatkan risiko kesalahan.

\textbf{Skenario Kode Asli Mentah:}
\begin{lstlisting}[style=pythonstyle]
def cek_validitas_pinjaman(a_kerja, b_jaminan, c_bebas_utang):
  
  # Validasi kualifikasi nasabah (versi ruwet):
  if not (not a_kerja and b_jaminan):
      # Cek tambahan di dalam:
      if a_kerja and (b_jaminan or not c_bebas_utang):
          return True
  return False
\end{lstlisting}

Representasi logikanya:
\[
\neg(\neg A \wedge B) \wedge (A \wedge (B \vee \neg C))
\]

\textbf{Langkah penyederhanaan:}
\begin{align*}
    1.\ & \neg(\neg A \wedge B) \wedge A \wedge (B \vee \neg C) & \text{\parbox[t]{5.8cm}{Bentuk awal}} \\
    2.\ & (A \vee \neg B) \wedge A \wedge (B \vee \neg C) & \text{\parbox[t]{5.8cm}{De Morgan}} \\
    3.\ & A \wedge (A \vee \neg B) \wedge (B \vee \neg C) & \text{\parbox[t]{5.8cm}{Komutatif}} \\
    4.\ & \mathbf{A} \wedge (B \vee \neg C) & \text{\parbox[t]{5.8cm}{Absorpsi}}
\end{align*}

\textbf{Transformasi Kode Bersih Pasca-Penyulingan (Refactoring):}
\begin{lstlisting}[style=pythonstyle]
def cek_validitas_pinjaman(a_kerja, b_jaminan, c_bebas_utang):
  
  # Filter Ekuivalensi Ceria nan Minimalis:
  return a_kerja and (b_jaminan or not c_bebas_utang)
\end{lstlisting}
Hasil refactoring menghasilkan logika yang lebih ringkas, lebih mudah dibaca, dan lebih mudah diuji.

\subsection{Aplikasi pada Rangkaian Gerbang Logika Digital}

Logika proposisional juga menjadi dasar optimasi rangkaian digital. Misalkan rancangan sinyal mengikuti ekspresi:
$$ Q_{output} = (A \land B) \lor (A \land \neg B) $$

\textbf{Penyederhanaan:}
\begin{align*}
    (A \wedge B) \vee (A \wedge \neg B) &\equiv A \wedge (B \vee \neg B) \quad \text{(Distributif)} \\
    &\equiv A \wedge \mathbf{T} \quad \text{(Hukum Negasi)} \\
    &\equiv \mathbf{A} \quad \text{(Identitas)}
\end{align*}

Dengan demikian, $Q_{output} \equiv A$. Artinya, rangkaian dapat disederhanakan secara signifikan karena keluaran cukup mengikuti sinyal $A$ \cite{wiki_proplogic}.


% ============================================================
% AKTIVITAS PEMBELAJARAN
% ============================================================

\begin{aktivitas}
  \item \textbf{Identifikasi Tautologi, Kontradiksi, Kontingensi:} Buat tabel kebenaran untuk $p \vee \neg p$, $p \wedge \neg p$, dan $p \rightarrow q$. Klasifikasikan masing-masing sebagai tautologi, kontradiksi, atau kontingensi.
  \item \textbf{Uji Ekuivalensi:} Buat tabel kebenaran untuk $\neg(p \wedge q)$ dan $\neg p \vee \neg q$. Buktikan bahwa keduanya ekuivalen.
  \item \textbf{Diskusi QA Sirkuit:} Berikan satu ekspresi logika (misalnya $A \wedge (A \vee B)$). Sederhanakan menggunakan hukum Absorpsi atau Dominasi. Bandingkan hasil dengan teman kelas.
  \item \textbf{Praktik Python:} Tulis fungsi \texttt{check\_equivalence(stmt1, stmt2)} yang memverifikasi $A \rightarrow B \equiv \neg A \vee B$ melalui tabel kebenaran 4 baris.
  \item \textbf{Prosedur Penyederhanaan:} Rangkum prosedur 4 langkah penyederhanaan (lucuti panah, De Morgan, Absorpsi, Distributif) dalam bagan alur vertikal.
  \item \textbf{Refactoring Kode:} Pilih potongan kode IF-ELSE dari program sederhana. Translasi ke logika proposisional, sederhanakan dengan hukum aljabar, tulis ulang kode yang bersih.
\end{aktivitas}

% ============================================================
% LATIHAN DAN REFLEKSI
% ============================================================

\begin{latihan}
  \item Buat tabel kebenaran untuk $p \vee \neg p$ dan $p \wedge \neg p$. Klasifikasikan masing-masing sebagai tautologi atau kontradiksi.
  \item Buat tabel kebenaran 8 baris (untuk $p, q, r$) untuk membuktikan Hukum Distributif: $p \vee (q \wedge r) \equiv (p \vee q) \wedge (p \vee r)$.
  \item Sederhanakan $\neg((\neg p \wedge q) \vee (p \vee \neg q))$ dengan hukum De Morgan, Asosiatif, dan Negasi Ganda (tanpa tabel).
  \item Buktikan $A \rightarrow (A \vee B)$ tautologi dengan mengubah ke bentuk $\neg A \vee (A \vee B)$ dan menerapkan hukum aljabar.
  \item Diberikan kode: \texttt{if (saldo\_cukup and (bukan\_kredit or bukan\_saldo\_cukup))}. Translasi ke variabel logika $s$ (saldo cukup) dan $k$ (kredit). Sederhanakan dan tulis refactoring kode yang bersih.
  \item \textbf{Refleksi:} Kapan Anda memilih tabel kebenaran dan kapan manipulasi aljabar untuk membuktikan ekuivalensi? Jelaskan kelebihan masing-masing.
\end{latihan}

% ============================================================
% ASESMEN
% ============================================================

\begin{asesmen}
\textbf{Instrumen Penilaian Evaluasi Kinerja (Sub-CPMK 1.2)}

\textbf{A. Pilihan Ganda}
\begin{enumerate}
  \item Hukum yang menyederhanakan $A \wedge (A \vee B)$ menjadi $A$ adalah:
  \begin{enumerate}
    \item Hukum De Morgan
    \item Hukum Idempoten
    \item Hukum Absorpsi
    \item Hukum Distributif
  \end{enumerate}
  Jawaban: (c)

  \item Proposisi $(p \vee \neg p) \rightarrow (q \wedge \neg q)$ bernilai kontradiksi karena sisi kiri tautologi dan sisi kanan kontradiksi, sehingga $(T \rightarrow F)$ bernilai:
  \begin{enumerate}
    \item Tautologi
    \item Kontingensi
    \item Disjungsi
    \item Kontradiksi
  \end{enumerate}
  Jawaban: (d)

  \item Hukum Komutatif berlaku untuk operasi yang sejenis, yaitu:
  \begin{enumerate}
    \item Implikasi
    \item Konjungsi dan disjungsi (AND atau OR)
    \item Negasi
    \item Biimplikasi
  \end{enumerate}
  Jawaban: (b)

  \item Pernyataan yang selalu benar disebut tautologi; selalu salah disebut kontradiksi. Pernyataan $p \rightarrow q$ termasuk:
  \begin{enumerate}
    \item Tautologi
    \item Kontradiksi
    \item Kontingensi
    \item Ekuivalensi
  \end{enumerate}
  Jawaban: (c)

  \item Ekuivalensi yang benar untuk implikasi $p \rightarrow q$ adalah:
  \begin{enumerate}
    \item $p \wedge \neg q$
    \item $\neg p \vee q$
    \item $\neg p \wedge q$
    \item $p \vee \neg q$
  \end{enumerate}
  Jawaban: (b)

  \item Menurut Hukum De Morgan, $\neg(p \wedge q)$ ekuivalen dengan:
  \begin{enumerate}
    \item $p \vee q$
    \item $\neg p \wedge \neg q$
    \item $\neg p \vee \neg q$
    \item $p \wedge q$
  \end{enumerate}
  Jawaban: (c)

  \item Bentuk $p \vee (q \wedge r) \equiv (p \vee q) \wedge (p \vee r)$ merupakan:
  \begin{enumerate}
    \item Hukum De Morgan
    \item Hukum Asosiatif
    \item Hukum Distributif
    \item Hukum Absorpsi
  \end{enumerate}
  Jawaban: (c)
\end{enumerate}

\textbf{B. Tugas Praktik}
\begin{itemize}
  \item Diberikan rumus $A \rightarrow (A \vee B)$. Ubah ke bentuk $\neg A \vee (A \vee B)$, terapkan hukum aljabar proposisi, dan tentukan apakah hasilnya tautologi, kontradiksi, atau kontingensi. Paparkan langkah penyelesaian dan justifikasi.
\end{itemize}

\textbf{Rubrik Penilaian}: Merujuk Lampiran Buku.
\end{asesmen}

% ============================================================
% CHECKLIST KOMPETENSI
% ============================================================

\begin{checklist}
  \item Saya mampu membedakan tautologi, kontradiksi, dan kontingensi.
  \item Saya mampu menggunakan tabel kebenaran untuk membuktikan ekuivalensi dua proposisi.
  \item Saya mampu menerapkan hukum Dominasi, Identitas, Negasi, Involusi, Idempoten, dan Absorpsi.
  \item Saya mampu menerapkan hukum Komutatif, Asosiatif, Distributif, dan De Morgan.
  \item Saya mampu mengubah implikasi $p \rightarrow q$ ke bentuk $\neg p \vee q$.
  \item Saya mampu menyederhanakan ekspresi dengan prosedur 4 langkah (panah, De Morgan, Absorpsi, Distributif).
  \item Saya mampu menerjemahkan kode IF-ELSE ke logika dan melakukan refactoring.
\end{checklist}

% ============================================================
% RANGKUMAN
% ============================================================

\begin{rangkuman}
Bab 4 membahas evaluasi proposisi majemuk melalui tabel kebenaran dan manipulasi aljabar logika untuk menentukan sifat serta ekuivalensinya.

\textbf{Poin-poin utama:}
\begin{itemize}
  \item Proposisi majemuk diklasifikasikan menjadi \textit{tautologi}, \textit{kontradiksi}, dan \textit{kontingensi} berdasarkan pola nilai kebenarannya.
  \item Ekuivalensi logis ($\equiv$) dapat dibuktikan dengan tabel kebenaran, yaitu ketika dua ekspresi memiliki kolom hasil yang identik pada seluruh kombinasi variabel.
  \item Hukum aljabar proposisi (dominasi, identitas, negasi, involusi, idempoten, absorpsi, komutatif, asosiatif, distributif, dan De Morgan) digunakan untuk menyederhanakan ekspresi secara sistematis.
  \item Implikasi dan biimplikasi dapat ditransformasikan ke bentuk operator dasar, misalnya $p \rightarrow q \equiv \neg p \vee q$.
  \item Strategi penyederhanaan bertahap membantu menghasilkan ekspresi yang lebih ringkas untuk analisis logika, refactoring kode, dan optimasi rangkaian digital.
\end{itemize}

\textbf{Kata kunci}: Tautologi, kontradiksi, kontingensi, ekuivalensi logis, tabel kebenaran, hukum aljabar proposisi, De Morgan, implikasi, biimplikasi.
\end{rangkuman}

\ifSubfilesClassLoaded{
  \renewcommand{\bibname}{Daftar Pustaka}
  \bibliographystyle{plain}
  \bibliography{../references}
}{}
\end{document}
